\mode*
\section{Related work}
\label{sec:related}

The Shell Tutor \autocite{Winder2024ShellTutor} is, to our knowledge, the
main teaching intervention aimed directly at the shell (and Git); its
logging infrastructure is also one of the few corpora of student shell
behaviour.
A second such corpus comes from cybersecurity training, where a
command-logging toolset records the shell commands students run and analyses
them for solution patterns and misconceptions
\autocite{Svabensky2021ShellCommands}.
The auto-graded exercise platform ShellOnYou is a further shell-teaching
intervention, reported particularly useful for cohorts with heterogeneous
Unix backgrounds \autocite{Berry2022ShellOnYou}; and the need reaches
beyond computer science---life-science curricula teach Unix command-line
skills using browser-embedded command lines precisely to remove
installation barriers \autocite{Khamvongsa2026UnixBiologists}.
Both collect behaviour; neither analyses it in variation-theoretic terms.
Our work is complementary: where the tutors scaffold \emph{doing}, we aim
to identify \emph{what must be discerned} and how to vary it.

Learners' understanding of the shell has, as far as we can find, been
\emph{measured} only rarely (\cref{app:protocol}).
The one phenomenographic study interviewed six students after a first
semester of Unix and found a hierarchical outcome space---Unix as a set
of commands, as a tool for solving certain problems, as a professional
computing environment---with indications that the level tracked the
final mark and the deep-learning score \autocite{DoyleLister2006Unix};
its teacher counterpart found academics' conceptions of Unix broadly
consistent with the students' \autocite{DoyleLister2007WhyUnix}.
No validated diagnostic instrument or concept inventory exists for the
shell or the command line; the closest is a pre- and post-course concept
inventory for operating systems, whose items concern the internals
(scheduling, paging, block maps) rather than the user's shell
\autocite{WebbTaylor2014OSCI}, and the computing concept inventories
reviewed by \textcite{Taylor2014CI} do not reach the shell either.
The nearest precedent in method reads students' conceptions of virtual
memory off written tests at the start and end of an operating-systems
course \autocite{Pamplona2024VirtualMemory}---the design we adopt for
the shell.
What the remaining empirical work measures is behaviour or
confidence: terminal histories matched against milestones
\autocite{Mirkovic2020Terminal}, auto-graded terminal tasks with a
confidence survey \autocite{BaileyZilles2019uAssign}, the
Carpentries' pre- and post-workshop surveys of self-reported skill
\autocite{Jordan2018Carpentry}, and self-efficacy surveys around
courses that teach the command line among other tools
\autocite{Gilson2022MissingCSClass}.
Our instrument (\cref{sec:instruments}) is, to our knowledge, the first
aspect-level diagnostic of the shell, and the first to pair it with
written accounts analysed phenomenographically.
Variation theory has been applied to computing education before, notably to
programming conceptions of non-CS engineering students
\autocite{ThuneEckerdal2009Variation} and to lab-session learning
\autocite{ThuneEckerdal2018Lab}; our companion papers
\autocite{VTProgMisconceptions,VTDebug} extend this line to misconception
catalogues and debugging.
The systematic search (\cref{app:protocol}) confirmed that the
learner-focused literature on the file system, on composition, and on the
remote shell remains sparse---the gaps that this paper's aspect analyses
address.
